\section{Muons}
\label{sec:ch6_muons}

\subsection{Reconstruction}

Muons are reconstructed using tracks in the ID and MS, together with their energy deposits in the calorimeters~\cite{ATLAS2021MuonID}.
The reconstructed muon types are:
\begin{itemize}
    \item \textbf{Combined (CB) muons} are reconstructed by matching and fitting ID and MS tracks, including the energy loss between the two tracking systems. An alternative method called the inside-out method is also used to increase the acceptance of low $\pT$ muons. In this method, MS hits are fitted together with the extrapolated ID track, without requiring an independently reconstructed MS track.    
    \item \textbf{Segment-tagged (ST) muons} are reconstructed by matching an ID track to one or more local track segments in the MS when a complete MS track is unavailable.
    \item \textbf{Calorimeter-tagged (CT) muons} are reconstructed by matching an ID track to calorimeter deposits compatible with a minimum-ionising particle. The muon acceptance is recovered in the central region $|\eta|<0.1$, where MS coverage is limited.
    \item \textbf{MS-extrapolated muons} are reconstructed by extrapolating an MS track towards the interaction region without a matched ID track.
\end{itemize}

\noindent CB muons are used as signal muons in this analysis, while ST and CT muons are additionally used as baseline muons.

\subsection{Identification}

Muon identification tests whether the detector hits and the ID and MS momentum measurements are consistent with a single muon trajectory.
This step is important because muon candidates can also be produced by in-flight decays of charged mesons, such as pions and kaons.
These decays can change the trajectory and momentum, leading to a disagreement between the momentum measurements.
For muon candidates, this agreement is quantified by the $q/p$ significance,
\begin{align*}
    S_{q/p}=\frac{\left|(q/p)_{\mathrm{ID}}-(q/p)_{\mathrm{MS}}\right|}
    {\sqrt{\sigma_{\mathrm{ID}}^2+\sigma_{\mathrm{MS}}^2}},
\end{align*}
where $(q/p)_{\mathrm{ID}}$ and $(q/p)_{\mathrm{MS}}$ are the charge-to-momentum ratios determined from the ID and MS tracks, respectively.
$\sigma_{\mathrm{ID}}$ and $\sigma_{\mathrm{MS}}$ denote the uncertainties.
Large values indicate poor agreement and help reject incorrectly matched candidates.

Muon identification provides \texttt{Loose}, \texttt{Medium} and \texttt{Tight} working points.
In this analysis, baseline muons are required to pass \texttt{Loose} identification while signal muons must additionally satisfy the \texttt{Medium} working point. The two working points are defined as follows:
\begin{itemize}
    \item \texttt{Medium}: CB muons within the ID acceptance must satisfy $S_{q/p}<7$. Outside $|\eta|<0.1$, they must also have at least two MS precision stations with at least three MDT or CSC hits each.
    \item \texttt{Loose}: the same requirements as \texttt{Medium}, with additional acceptance for ST and CT muons in $|\eta|<0.1$, where MS coverage is limited.
    It also accepts additional inside-out CB muons with $\pT<7~\GeV$.
\end{itemize}
The typical prompt-muon reconstruction and identification efficiencies are about 99\% for \texttt{Loose} and 97\% for \texttt{Medium}.
The corresponding representative rejection rates against light hadrons are about 99.75\% and 99.83\%, respectively~\cite{ATLAS2021MuonID}.

\subsection{Isolation}

Muons from heavy-flavour decays are often accompanied by other particles inside a jet.
They are suppressed using track and calorimeter isolation, calculated similarly to electron isolation, with the muon's own contribution excluded.
Baseline muons are required to pass \texttt{Loose\_VarRad} isolation, and signal muons must additionally satisfy \texttt{Tight\_VarRad}.
Both working points use a $\pT$-dependent cone radius for the track isolation sum and a fixed cone of $\Delta R=0.2$ for the calorimeter isolation sum. The specific cut thresholds are defined as follows:
\begin{itemize}
    \item \texttt{Loose\_VarRad}: the track and calorimeter isolation sums must be below 15\% and 30\% of the muon $\pT$, respectively.
    \item \texttt{Tight\_VarRad}: the track and calorimeter isolation sums must be below 4\% and 15\% of the muon $\pT$, respectively.
\end{itemize}
The typical prompt-muon isolation efficiencies are about 95\% for \texttt{Loose\_VarRad} and 86\% for \texttt{Tight\_VarRad}.
The corresponding representative rejection rates against non-prompt muons are about 95\% and 98\%, respectively~\cite{Maricic2025MuonIsolation}.

\subsection{Calibration}

The calibration for muons follows a procedure similar to that for electrons and consists of a momentum calibration and an efficiency calibration.

For the momentum calibration, momentum biases in the track reconstruction caused by detector misalignment are first corrected in data.
The momentum scale and resolution are then corrected for simulated events using the $m_{\mu\mu}$ distributions in $\zmumu$ and $J/\psi\rightarrow\mu\mu$ events~\cite{ATLAS2023MuonCalibration}.

For the efficiency calibration, the reconstruction, identification and isolation efficiencies are measured with the tag-and-probe method using $\zmumu$ events.
$J/\psi\rightarrow\mu\mu$ events are also used to measure the reconstruction and identification efficiencies at low $\pT$~\cite{ATLAS2021MuonID}.
The scale factors defined in Eq.~\eqref{eq:ch6_efficiency_sf} are then applied to simulated event weights for the muon working points used in this analysis to correct the muon selection efficiencies.
